\begin{proofbox}
\textbf{推导（只用三角参数法）。}

椭圆动点设为
\[
P(\theta)=(a\cos\theta,b\sin\theta).
\]
由两点距离公式，
\[
F(\theta)=(a\cos\theta-x_0)^2
          +(b\sin\theta-y_0)^2.
\]
由于 $F$ 连续且以 $2\pi$ 为周期，
它在一整周上能取到最大值和最小值。

\textbf{第一步：求导。}
\[
\begin{aligned}
F'(\theta)
={}&-2a(a\cos\theta-x_0)\sin\theta\\
   &+2b(b\sin\theta-y_0)\cos\theta.
\end{aligned}
\]
令 $F'(\theta)=0$，整理得
\[
(a^2-b^2)\sin\theta\cos\theta
-ax_0\sin\theta+by_0\cos\theta=0.\tag{1}
\]

\textbf{第二步：统一消去三角函数。}
令 $t=\tan\frac{\theta}{2}$（暂不取 $\theta=\pi$），
则
\[
\sin\theta=\frac{2t}{1+t^2},\qquad
\cos\theta=\frac{1-t^2}{1+t^2}.
\]
把它们代入 (1)，并乘以 $(1+t^2)^2$：
\[
\begin{aligned}
0={}&2(a^2-b^2)t(1-t^2)\\
   &-2ax_0t(1+t^2)\\
   &+by_0(1-t^2)(1+t^2).
\end{aligned}
\]
展开、整理后得到关于 $t$ 的四次以内方程：
\[
\boxed{\begin{aligned}
0={}&by_0t^4+2(a^2-b^2+ax_0)t^3\\
   &+2(ax_0-a^2+b^2)t-by_0.
\end{aligned}}
\tag{2}
\]

\textbf{第三步：还原点并比较。}
对 (2) 的每个\textbf{实根} $t$，
还原椭圆上的候选点
\[
P_t=\left(a\frac{1-t^2}{1+t^2},
            b\frac{2t}{1+t^2}\right).
\]
再\textbf{单独检查}万能代换未覆盖的
$\theta=\pi$，即 $P_\pi=(-a,0)$。
比较所有候选点的
\[
MP^2=(x-x_0)^2+(y-y_0)^2
\]
的最小值和最大值，最后分别开方。

由一元可导周期函数的极值必要条件，
全局距离最值必在这些候选点中取得，故方法成立。
\end{proofbox}
